\documentclass{book}

\usepackage{fontspec} % used to import Calibri
\usepackage{anyfontsize} % used to adjust font size

% needed for inch and other length measurements
% to be recognized
\usepackage{calc}

% for colors and text effects as is hopefully obvious
\usepackage[dvipsnames]{xcolor}
\usepackage{soul}

% control over margins
\usepackage[margin=1in]{geometry}
\usepackage[strict]{changepage}

\usepackage{mathtools}
\usepackage{amsfonts}
\usepackage{bm}

\usepackage[scr=rsfso, scrscaled=.96]{mathalpha}

% This is how I'm getting the nice caligraphy font :(
\DeclareMathAlphabet{\eulerscr}{U}{eus}{m}{n}
\newcommand{\mathcalli}[1]{\text{\scalebox{1.11}{$\eulerscr{#1}$}}}


\usepackage{amssymb} % originally imported to get the proof square
\usepackage{xfrac}
\usepackage[overcommands]{overarrows} % Get my preferred vector arrows...
\usepackage{relsize}

\usepackage{tikz}
   \usetikzlibrary{arrows.meta}
   \usetikzlibrary{graphs, graphs.standard}

\usepackage{quiver} %commutative diagrams




\newfontfamily{\calibri}{Calibri}
\setlength{\parindent}{0pt}
\definecolor{RawerSienna}{HTML}{945D27}

% ~~~~~~~~~~~~~~~~~~~~~~~~~~~~~~~~~~~~~~~~~~~~~~~~~~
%Arrow Commands:

% Thank you Bernard, gernot, and Sigur who I copied this from:
% https://tex.stackexchange.com/questions/364096/command-for-longhookrightarrow
\renewcommand{\hookrightarrow}{\lhook\joinrel\rightarrow}
\renewcommand{\hookleftarrow}{\leftarrow\joinrel\rhook}
\newcommand{\hooklongrightarrow}{\lhook\joinrel\longrightarrow}
\newcommand{\hooklongleftarrow}{\longleftarrow\joinrel\rhook}
\newcommand{\hookxlongrightarrow}[2][]{\lhook\joinrel\xrightarrow[#1]{#2}}
\newcommand{\hookxlongleftarrow}[2][]{\xleftarrow[#1]{#2}\joinrel\rhook}

% Thank you egreg who I copied from:
% https://tex.stackexchange.com/questions/260554/two-headed-version-of-xrightarrow
\newcommand{\longrightarrowdbl}{\longrightarrow\mathrel{\mkern-14mu}\rightarrow}
\newcommand{\rightarrowdbl}{\rightarrow\mathrel{\mkern-14mu}\rightarrow}
\newcommand{\longleftarrowdbl}{\leftarrow\mathrel{\mkern-14mu}\longleftarrow}

\newcommand{\xrightarrowdbl}[2][]{%
  \xrightarrow[#1]{#2}\mathrel{\mkern-14mu}\rightarrow
}
\newcommand{\xleftarrowdbl}[2][]{%
  \leftarrow\mathrel{\mkern-14mu}\xleftarrow[#1]{#2}
}

\newcommand{\locUConverges}{%
	\xrightarrow{\ell.u.}
}
\newcommand{\UConverges}{%
	\xrightarrow{u.}
}



% ~~~~~~~~~~~~~~~~~~~~~~~~~~~~~~~~~~~~~~~~~~~~~~~~~~

\newcommand{\hTwo}{%
\color{Black}%
   \fontsize{13}{15}\selectfont%
}
% \newcommand{\scratchWork}{%
%    \color{PineGreen!85!Orange}
%    \fontsize{12}{14}\selectfont%
% }
\newcommand{\hThree}{%
   \color{Black}%
   \fontsize{12}{14}\selectfont%
}
\newcommand{\myComment}{%
   \color{RawerSienna}%
   \fontsize{12}{14}\selectfont%
}
\newcommand{\pracOne}{
   \color{BrickRed}%
   \fontsize{13}{15}\selectfont%
}
\newcommand{\pracTwo}{
   \color{Orange}%
   \fontsize{12}{14}\selectfont%
}
\newcommand{\exTwo}{%
   \color{Purple}%
   \fontsize{13}{15}\selectfont%
}
\newcommand{\exThree}{%
   \color{Purple}%
   \fontsize{12}{14}\selectfont%
}
\newcommand{\exThreeP}{%
   \color{RedViolet}%
   \fontsize{12}{14}\selectfont%
}
\newcommand{\exFourP}{%
   \color{RedViolet}%
   \fontsize{11}{13}\selectfont%
}
\newcommand{\exPPP}{%
   \color{VioletRed}%
   \fontsize{12}{14}\selectfont%
}

% Homework standard below (God the bloat in the header is absurd...)
% ~~~~~~~~~~~~~~~~~~~~~~~~~~~~~~~~~~~~~~~~~~~~~~~~
\newcommand{\Hstatement}{%
   \color{MidnightBlue!90!Black}%
   \fontsize{12}{13}\selectfont%
}
\newcommand{\HexOne}{%
   \color{Purple}%
   \fontsize{12}{13}\selectfont%
}
\newcommand{\HexTwoP}{%
   \color{RedViolet}%
   \fontsize{12}{13}\selectfont%
}
\newcommand{\HexPPP}{%
   \color{VioletRed}%
   \fontsize{11}{12}\selectfont%
}

% ~~~~~~~~~~~~~~~~~~~~~~~~~~~~~~~~~~~~~~~~~~~~~~~~

\newcommand{\cyPen}[1]{{\vphantom{.}\color{Cerulean}#1}}
\newcommand{\redPen}[1]{{\vphantom{.}\color{Red}#1}}

\newenvironment{myIndent}{%
   \begin{adjustwidth}{2.5em}{0em}%
}{%
   \end{adjustwidth}%
}

\newenvironment{myDindent}{%
   \begin{adjustwidth}{5em}{0em}%
}{%
   \end{adjustwidth}%
}

\newenvironment{myTindent}{%
   \begin{adjustwidth}{7.5em}{0em}%
}{%
   \end{adjustwidth}%
}

\newcommand{\udefine}[1]{{%
   \setulcolor{Red}%
   \setul{0.14em}{0.07em}%
   \ul{#1}%
}}

\newcommand{\blab}[1]{\textbf{#1}}
\newcommand{\blect}[1]{{\color{MidnightBlue}\textbf{#1}}}

\newcommand{\uuline}[2][.]{%
{\vphantom{a}\color{#1}%
\rlap{\rule[-0.18em]{\widthof{#2}}{0.06em}}%
\rlap{\rule[-0.32em]{\widthof{#2}}{0.06em}}}%
#2}

\newcommand{\pprime}{{\prime\prime}}
\newcommand{\suchthat}{ \hspace{0.3em}s.t.\hspace{0.3em}}
\newcommand{\rea}[1]{\mathrm{Re}(#1)}
\newcommand{\ima}[1]{\mathrm{Im}(#1)}
\newcommand{\comp}{\mathsf{C}}
\newcommand{\trans}{\mathsf{T}}
\newcommand{\myHS}{ \hspace{0.5em}}
\newcommand{\gap}{\phantom{2}}

\newcommand{\GenLin}{\ensuremath{\mathrm{GL}}}
\newcommand{\Cay}{\ensuremath{\mathrm{Cay}}}

\newcommand{\myId}{\mathrm{Id}}
\newcommand{\myIm}{\mathrm{im}}
\newcommand{\Obj}{\mathrm{Obj}}
\newcommand{\Hom}{\mathrm{Hom}}
\newcommand{\End}{\mathrm{End}}
\newcommand{\Aut}{\mathrm{Aut}}

\newcommand{\df}{\mathrm{d}}
\newcommand{\Df}{\mathrm{D}}

\newcommand{\mcateg}[1]{{\bm{\mathsf{#1}}}}

\newcommand{\mdeg}{\mathrm{mdeg}\phantom{.}}

\newcommand{\dividesDeprecated}{\mathop{\mid}}
\newcommand{\divides}{\mathrel{\mid}}

\newcommand{\card}{\mathrm{card}}
\newcommand{\supp}{\mathrm{supp}}
\newcommand{\diam}{\mathrm{diam}}
\newcommand{\conv}{\mathrm{conv}}
\newcommand{\opnorm}{\mathrm{op}}
\newcommand{\loc}{\mathrm{loc}}
\newcommand{\sgn}{\mathrm{sgn}}
\newcommand{\acc}{\mathrm{acc}}
\newcommand{\mSpan}{\mathrm{span}}
\newcommand{\Interior}{\mathop{\mathrm{Int}}}

\newcommand{\mMat}[1]{\mathbf{#1}}

\newcommand{\NBV}{\ensuremath{\mathrm{NBV}}}
\newcommand{\Acc}{\mathrm{Acc}}
\newcommand{\BV}{\ensuremath{\mathrm{BV}}}
\newcommand{\Var}{\ensuremath{\mathrm{Var}}}

\newcommand{\Alt}{\mathrm{Alt}}
\newcommand{\Sym}{\mathrm{Sym}}

\newcommand{\weakst}{weak$^*$ }

\newcommand{\radtimes}{\mathop{\widehat{\times}}}

\newcommand{\mMod}[1]{\phantom{a}(\mathrel{\mathrm{mod}} #1)}
\newcommand{\Fun}{\mathrm{Fun}}
\newcommand{\act}{\mathrm{act}}
\newcommand{\Fix}{\mathrm{Fix}}
\newcommand{\Sub}{\mathrm{Sub}}
\newcommand{\Cl}{\mathrm{Cl}}
\newcommand{\GL}{\mathrm{GL}}
\newcommand{\SL}{\mathrm{SL}}
\newcommand{\PSL}{\mathrm{PSL}}
\newcommand{\core}{\mathrm{core}}
\newcommand{\Syl}{\mathrm{Syl}}
\newcommand{\Iso}{\mathrm{Iso}}
\newcommand{\Homeo}{\mathrm{Homeo}}
\newcommand{\Inn}{\mathrm{Inn}}
\newcommand{\Out}{\mathrm{Out}}
\newcommand{\ab}{\mathrm{ab}}
\newcommand{\Max}{\mathrm{Max}}
\newcommand{\lt}{\mathrm{lt}}
\newcommand{\Nil}{\mathrm{Nil}}
\newcommand{\Ideal}{\mathrm{Ideal}}
\newcommand{\Spec}{\mathrm{Spec}}
\newcommand{\Res}{\mathrm{Res}}
\newcommand{\prim}{\mathrm{prim}}
\newcommand{\exConv}{\mathrm{ex}}
\newcommand{\mRank}{\mathrm{rank}}
\newcommand{\DivGroup}{\mathrm{Div}}
\newcommand{\divGroup}{\mathrm{div}}
\newcommand{\PrinGroup}{\mathrm{Prin}}
\newcommand{\Ann}{\mathrm{Ann}}
\newcommand{\Tor}{\mathrm{Tor}}
\newcommand{\diag}{\mathrm{diag}}
\newcommand{\Nat}{\mathrm{Nat}}
\newcommand{\Embed}{\mathrm{Embed}}
\newcommand{\Isom}{\mathrm{Isom}}
\newcommand{\mChar}{\mathrm{char}}
\newcommand{\Gal}{\mathrm{Gal}}
\newcommand{\Int}{\mathrm{Int}}
\newcommand{\myTr}{\mathrm{Tr}}
\newcommand{\genus}{\mathrm{genus}}
\newcommand{\Ass}{\mathrm{Ass}}
\newcommand{\Sin}{\mathrm{Sin}}

\DeclareMathOperator{\lcm}{lcm}
\DeclareMathOperator{\Log}{Log}
\DeclareMathOperator{\ord}{ord}
\DeclareMathOperator{\adj}{adj}
\DeclareMathOperator{\symdif}{\triangle}
\DeclareMathOperator{\Average}{Average}
\DeclareMathOperator*{\AverageAst}{Average}

\newcommand{\th}[1]{\ifnum\value{#1}}



% Thank you Gonzalo Medina and Moriambar who wrote this on stack exchange:
%https://tex.stackexchange.com/questions/74125/how-do-i-put-text-over-symbols%
\newcommand{\myequiv}[1]{\stackrel{\mathclap{\mbox{\footnotesize{$#1$}}}}{\equiv}}

% Thank you chs who wrote this on stack exchange:
%https://tex.stackexchange.com/questions/89821/how-to-draw-a-solid-colored-circle%
\newcommand{\filledcirc}[1][.]{\ensuremath{\hspace{0.05em}{\color{#1}\bullet}\mathllap{\circ}\hspace{0.05em}}}

%Thank you blerbl who wrote this on stack exchange:
%https://tex.stackexchange.com/questions/25348/latex-symbol-for-does-not-divide
\newcommand{\ndiv}{\hspace{-0.3em}\not|\hspace{0.35em}}

\newcommand{\mySepTwo}[1][.]{%
   {\noindent\color{#1}{\rule{6.5in}{0.5mm}}}\\%
}
\newcommand{\mySepThree}[1][.]{%
   {\noindent\color{#1}{\rule{6in}{0.25mm}}}\\%
}
\newcommand{\mySepFour}[1][.]{%
   {\noindent\color{#1}{\rule{5.5in}{0.12mm}}}\\%
}

\newcommand{\retTwo}{\hfill\bigbreak}

% Hopefully this will be good enough
\NewOverArrowCommand{myVector}{%
   start = {{\smallermathstyle\relbar}},
   middle = {{\smallermathstyle\relbar}},
   end={{\rightharpoonup}}, space before arrow=0.15em,
   space after arrow=-0.045em,
}
\newcommand{\mVec}[1]{\myVector{#1}}
\newcommand{\mVecAst}[1]{\myVector*{#1}}

\begin{document}
\setul{0.14em}{0.07em}
\calibri

\exTwo\ul{Lemma 1.4:} Let $M$ be a paracompact Hausdorff space and $\{U_\alpha\}_{\alpha \in A}$ \textit{any} open cover. Then there exists a locally finite open cover $\{V_\alpha\}_{\alpha \in A}$ such that $\overline{V_\alpha} \subseteq U_\alpha$ for all $\alpha$.
\begin{myIndent}\exThreeP
	Proof:\\
	Our first step is showing that paracompact and Hausdorff implies regular (i.e. $T_3$).
	\begin{myIndent}\exPPP
		Let $A \subseteq M$ be a closed set and suppose $p \notin A$. Then since $M$ is Hausdorff, for all $q \in A$ we can pick an open set $V_q \subseteq M$ such that $p \notin \overline{V_q}$. Next, note that $\{M - A\} \cup \{V_q\}_{q \in A}$ is an open cover of $M$. So by paracompactness, we can take a locally finite refinement $\{W_\beta\}_{\beta \in B}$ covering $M$.\retTwo

		Let $B^\prime$ be the collection of $\beta \in B$ such that $W_\beta \cap A \neq \emptyset$. Since we know that\\ $M = \bigcup_{\beta \in B} W_\beta$, we can thus conclude that $A \subseteq \bigcup_{\beta \in B^\prime} W_\beta$. At the same time though, because $\{W_\beta\}_{\beta \in B^\prime}$ is locally finite, we can say that:
		
		{\centering$\overline{\bigcup_{\beta \in B^\prime}W_\beta} = \bigcup_{\beta \in B^\prime} \overline{W_\beta} \subseteq \bigcup_{q \in A} \overline{V_q}$ where the latter set doesn't contain $p$.\retTwo\par}

		So, the open sets $U \coloneqq \bigcup_{\beta \in B^\prime} W_\beta$ and $V \coloneqq M - \overline{U}$ separate $p$ and the closed set $A$.\retTwo
	\end{myIndent}

	Next, let $\mathcalli{A}$ be the collection of open subsets of $M$ whose closures are contained in a set $U_\alpha$. Because $M$ is regular, we know that $\mathcalli{A}$ covers all of $M$. So by paracompactness, we can take a locally finite refinement $\{W_\beta\}_{\beta \in B}$ of $\mathcalli{A}$ covering $M$.\retTwo

	Define a function $f : B \to A$ such that $\overline{W_\beta} \subseteq U_{f(\beta)}$ for all $\beta$. Then for every $\alpha$ define $V_\alpha \coloneqq \bigcup_{\beta \in f^{-1}(\{\alpha\})} W_\beta$. Clearly $\bigcup_{\alpha \in A} V_\alpha = \bigcup_{\beta \in B} W_\beta = M$. Also note that because the $W_\beta$ are locally finite:

	{\centering $\overline{V_\alpha} = \overline{\bigcup_{\beta \in f^{-1}(\{\alpha\})} W_\beta} = \bigcup_{\beta \in f^{-1}(\{\alpha\})} \overline{W_\beta} \subseteq U_\alpha$. $\blacksquare$\retTwo\par}
\end{myIndent}

\ul{Theorem 2.23:} Let $M$ be a paracompact Hausdorff space and let $\mathcalli{U} = \{U_\alpha\}_{\alpha \in A}$ be\\ \textit{any} open cover. Then there is a continuous partition of unity $\{f_\alpha: M \to [0, 1]\}_{\alpha \in A}$\\ subordinate to $\mathcalli{U}$ such that $\supp(f_\alpha) \subseteq U_\alpha$ for all $\alpha \in A$.
\begin{myIndent}\exThreeP
	Proof:\\
	We know from before that paracompact and Hausorff implies regular. As it turns out, we can also use an identical argument (except that we replace $p$ with a disjoint closed set $B$ from $A$) in order to prove that paracompact and regular implies normal (i.e. $T_4$). This is useful because we can apply Urysohn's lemma to any normal space.\retTwo

	By lemma 1.4, we can find two locally finite open covers $\{V_\alpha\}_{\alpha \in A}$ and $\{W_\alpha\}_{\alpha \in A}$ such that $\overline{W_\alpha} \subseteq V_\alpha \subseteq \overline{V_\alpha} \subseteq U_\alpha$ for all $\alpha$. And then by Urysohn's lemma we can for every $\alpha \in A$ find a function $g \in C(M, [0, 1])$ such that $g \equiv 1$ on $\overline{W_\alpha}$ and $g \equiv 0$ on $M - V_\alpha$.\retTwo

	But now note that as $\supp(g_\alpha) \subseteq \overline{V_\alpha}$ and $\{\overline{V_\alpha}\}_{\alpha \in A}$ is locally finite, we know that every $p \in M$ has a neighborhood on which only finitely many $g_\alpha$ are nonzero. So, we can define a continuous function $h = \sum_{\alpha \in A} g_\alpha$. And then we define $f_\alpha = g_\alpha / h$ for all $\alpha$. This satisfies what we want. $\blacksquare$\newpage
\end{myIndent}

\pracOne Here is a proposition about locally finite collections of sets.\retTwo

\ul{Exercise 1.14:} Suppose $\mathcalli{X}$ is a locally finite collection of subsets of a topological space $M$.
\begin{itemize}
	\item[(a)] The collection $\{\overline{X} : X \in \mathcalli{X}\}$ is also locally finite.
	\begin{myIndent}\pracTwo
		Proof:\\
		Let $x \in M$ and let $U \subseteq M$ be an open set containing $x$ that intersects only the\\ elements of a finite subset $\mathcalli{Y}$ of $\mathcalli{X}$. Now we want to show that if $X \in \mathcalli{X}$ satisfies\\ $\overline{X} \cap U \neq \emptyset$, then $X \in \mathcalli{Y}$. Luckily, if $\overline{X} \cap U \neq \emptyset$ but $X \cap U = \emptyset$ so that $X \notin \mathcalli{Y}$, then it must be that any $x \in \overline{X} \cap U$ is an accumulation point of $X$. However, that is immediately contradicted by the fact that $U$ is a neighborhood of $x$ which does not intercept $X$ anywhere. $\blacksquare$\retTwo	
	\end{myIndent}
	
	\item[(b)] $\overline{\bigcup_{X \in \mathcalli{X}} X} = \bigcup_{X \in \mathcalli{X}} \overline{X}$.
	
	\begin{myIndent}\pracTwo
		Proof:\\
		It is always true that $\bigcup_{X \in \mathcalli{X}} \overline{X} \subseteq \overline{\bigcup_{X \in \mathcalli{X}} X}$.\retTwo

		Meanwhile, to show the other inclusion, suppose $x$ is an accumulation point of\\ $\bigcup_{X \in \mathcalli{X}} X$ but not in any individual $\overline{X}$. Then it must be the case that every\\ neighborhood of $x$ intercepts some set in $\mathcalli{X}$, yet it must also be the case that for every $X \in \mathcalli{X}$ there is a neighborhood of $x$ with doesn't intercept $X$. The only way this is possible is if every neighborhood intercepts infinitely many $X \in \mathcalli{X}$.
		\begin{myIndent}
			Let $N$ be any neighborhood of $x$. Then we can construct an infinite sequence of distinct sets $(X_n)_{n \in \mathbb{N}}$ in $\mathcalli{X}$ intercepting $N$ as follows.\retTwo

			For the ease of notation set $U_0 = M$. Now at step $n$, choose any $X_n \in \mathcalli{X}$ that intersects $N \cap \bigcap_{k=0}^{n-1} U_k$. Note that we can do this since a finite intersection of neighborhoods of $x$ is still a neighborhood of $x$. But now we can choose some neighborhood $U_{n}$ of $x$ such that $X_n \cap U_n = \emptyset$. And now repeat this reasoning.\retTwo

			It's clear that all the chosen $X_n$ intercept $\bigcap_{k=0}^{n-1} U_k \subseteq N$. Additionally, all the chosen $X_n$ are distinct since $X_n \cap \bigcap_{k=0}^{N} U_k \subseteq X_n \cap U_n = \emptyset$ for all $N \geq n$.\retTwo
		\end{myIndent}

		But that contradicts that $\mathcalli{X}$ is locally finite. $\blacksquare$\retTwo
	\end{myIndent}
\end{itemize}

\end{document}